\documentclass[preprint,12pt, a4paper]{elsarticle}

\usepackage{amsmath}
\usepackage{amsfonts}
\usepackage{graphicx}
\usepackage{subcaption}


\usepackage[colorlinks=true,hyperfootnotes=true,breaklinks=true]{hyperref}
\usepackage[capitalise]{cleveref}

	

\begin{document}

\title{Nonparametric comparison of  input output curves in electrophysiology.} 

\author[label1]{Patrick Michalik}
\author[label2]{Tomasz Wójtowicz}
\author[label2]{Jakub Włodarczyk}
\author[label3]{Błażej Ruszczycki}
\address[label1]{AGH University of Science and Technology, Faculty of Electrical Engineering,
Automatics, Computer Science and Biomedical Engineering, al. Mickiewicza 30 B1, 30-059 Kraków, Poland}
\address[label2]{Nencki Institute of Experimental Biology, Pasteura 3, 02-093 Warsaw, Poland}
\address[label3]{AGH University of Science and Technology, Faculty of Physics and Applied Computer Science, al. Mickiewicza 30 D10, 30-059 Kraków, Poland}

\begin{abstract}
    o
    
\end{abstract}

\date{\today}

\maketitle

\section{Introduction}

 Nonparametric tests relax distribution assumptions \cite{krzywinski}
comparison of regression curves \cite{Neumeyer}    


\section{Nonparametric statistics}

In the following section, we present the computational approach, used to estimate the p-value, using nonparametric randomization.
Initially, for the reason of simplicity, we consider only two experimental groups. Later, the method is generealized to arbitrary number of groups.

In the presented design, the basic object with which we operate, is the input-output curve, denoted as $f_{g,s,m}(x)$, where $g$ denotes the experimental group, $s$ denotes the subject in the experimental group, $m$ denotes the number of the measrurement for the given subject and $x$ denotes the contiuous within-subject parameter (the excitation voltage).  For the null hypothesis $H_0$,we assume that the input-output curves do not depend statistically on the experimental groups.  The null hypothesis $H_0$ is represented by the set of curves in the  randomized ensemble $\tilde{f}^{[k]}_{g,a,s}(x)$, where $k$ numbers the subsequent randomizations.

The p-value is computed by performing the ranking of the difference between the actual curves and the curves from the randomized ensemble.
In order to prepare such a ranking, we need to compare quantitatively different input-output curves.

\section{Results}
\section{Materials and methods}

\subsection{Randomized ensemble}
We perform the conditional randomization, creating new expermental gru
\subsection{L2 measure}

\section{Discussion}

 \bibliography{inout} 
\bibliographystyle{ieeetr}

 
\end{document}
